\subsection{Ground Truth Dataset Collection}
Each datapoint consists of an image, a question, and measured value. The categories we consider are:

\textbf{Weight} - When posing questions about weight, we avoid asking about standard packaging. For instance, instead of asking
about the weight of a packaged quantity of sugar, that's likely to be in standard weights, we instead draw some amount of sugar
and ask about its weight. To do well on such questions, one must be familiar with (i) relative densities of materials, (ii) be able to infer the size or quantity of an object by spatial reasoning from other objects in the scene. We used a kitchen scale, 
a luggage scale (with hook), and a bathroom weighing scale to obtain the ground truth, depending on size of object.

\textbf{Volume} - This consists of two types: (i) how much volume of liquid a container is holding or can hold, 
(ii) what is the volume of an object itself. We used beakers of various capacities to measure volume. For (ii), we submerged the object in water and measured volume of displaced liquid. Since the objects often have complex geometries, spatial reasoning by itself is insufficient to do well on these questions.

\textbf{Angle} - We ask about the angle of inclination between two objects. The angle to be determined is annotated in image with a hand-drawn arc. We measure the angle with either a protractor or a digital inclinometer, depending on the setting.

\textbf{Degree of stability} - In the image, we show a stable configuration of objects and ask about the degree of perturbation needed to make it unstable. The choice of this type of question is motivated by two considerations: (i) the answer is a numerical value, (ii) it is hard to capture a configuration in a unstable state; it's much more likely that the scene shown in the image will correspond to a stable state. This makes the binary version of the question easy to answer. To get the ground truth, we perturb the object in question and measure the point at which the configuration of the scene becomes unstable.

\textbf{Geometric Fit} - We ask here whether one object has enough vacant space to fully contain another object. Often the object to be contained is deformable and our question permits bending, twisting such objects to the extend that the object is not damaged. Had both the objects been non-deformable, our question would involve just spatial reasoning - estimation of dimensions of the two objects. To do well in our setting, however, it helps to have a feel of how pliable objects are and the degrees to which they can be deformed. The ground truth is obtained by us trying to make the object in question fit within the other and noting the result.

Our considerations when collecting the dataset were:

\begin{itemize}
    \item The values should span a wide range. We've attempted to ask questions about small, medium and large objects in the scene. This was more feasible for weight and angle. Due to difficulties in capturing ground truth, we mostly considered small and medium objects for volume, degree of stability, and geometric fit questions.
    \item The questions test for abilities beyond spatial reasoning. Strong priors, such as for relative density or flexibility, are needed to answer the questions correctly that the image
    is often insufficient to convey fully and gets instilled by physical interactions. We note that spatial reasoning is also very important to do well on our tests. This is strongly the case for angle and to some extent volume estimation.
    \item To the extent possible, we've tried to capture images outdoors. This brings more heterogeneity - for instance with natural lighting and texture variations.
    \item We've ensured that there are plenty of reference objects around in the scene and clearly visible. So for instance, we avoid images where a single image is captured in close-up. The reference objects provide anchors and make it more feasible to answer physical property questions about another object. In our daily experience, humans also draw upon such references and comparisons.
    \item We've tried to capture objects from atypical angles to the extent that it's feasible to answer the question posed. This means that often lateral viewpoints were preferred over frontal viewpoints. This is to make the estimation more difficult and draw more strongly upon priors.
    \item The final basis for deciding whether to include a datapoint in the dataset is whether there is enough information in the image to make it possible for humans to take a meaningful shot at answering the question. We ensure that there are no trick questions or setups that cause visual illusions.
\end{itemize}

\subsection{Human Baseline}
To evaluate how well humans do on these estimation tasks, we created a web application whereby on each trial a user is presented with a subset of datapoints to answer. We ensure that each user sees novel samples during each trial (until all samples in the bechmark dataset are exhausted). We conduct this trial with multiple individuals to capture a sense of how well humans generally do on these tasks. To make it interesting for participants, we presented the trials as a game whereby the participant can compare their performance to frontier models. To this end, we made use of LLM responses collected as detailed in the following section. 

\subsection{Simple Evaluation of Frontier LLMs}
  We evaluate closed weight frontier models on our benchmark dataset. For each family of model, we select the best (largest) model with reasoning effort set to High. We present the image and question to the frontier models and collect the responses. The response consists of both the reasoning trace and the final predicted answer. We use structured output so as to avoid having to parse the predictions. We collect multiple responses from each model by varying the seed. 

For categories with numerical answers such as weight, volume, angle, and degree of stability, we report average percentage error over all samples in that category. Geometric fit is a binary classification and we report precision and recall for it. For each category, we also report the standard deviation over multiple runs for each datapoint, averaged over all datapoints.

For the reasoning traces, we perform qualitative analysis. [TODO]

\subsection{Robustness Evaluation of Frontier LLMs}
  Consider the images for our estimation questions. If we blur them a little, a human being is still likely to reason tFor the kind of estimation tasks we are interested in, it suffices to know the basic shapes, identities and orientations of objects in the image. Some kinds of information like brand names of products can make the VLMs fall back on recall from memory. To strongly test for how well the VLMs are able to reason, we introduce varying levels of blur - low, medium, and high - in each image. We implement this using Gaussian kernel with varying standard deviations. Note that the important information to be able to answer 


\subsection{Interventions to Improve Model Performance}

We aim to improve the peformance of VLMs on our estimation tasks by the following interventions. We tried prompt template and fine-tuning work on small and medium open-weight VLMs. The tool augmentation approach was used for closed weight frontier models.

\subsubsection{Intervention 1: Prompt Template}
We investigate whether giving a detailed system prompt can help guide a model to make more accurate inferences. To achieve this, we use a form of model distillation we describe now. We first do benchmark evaluation on the target smaller models. For each model and for category, we identify the samples where the target model does noticeably worse than a good quality closed model. For each category, the system prompt is derived by comparing the reasoning trace of the small target model and larger closed model for a set of datapoints in that category. Specifically, given these pairs, we ask a large closed model to identify the differences in patterns of reasoning between the well-performing larger model and the worse small model. This difference generalizes over the pairs in the entire set for that category.

We evaluate the small model again on the benchmark dataset but this time it is given the prompt template we derived as system prompt.

\subsubsection{Intervention 2: Tool Augmentation}
In this intervention, we aim to improve the performance of a small VLM by giving it more contextual information for each datapoint to help it reason better. The contextual information is derived by running external tools offline for each datapoint and then giving that information in a structured manner as part of the input prompt of questioning. These were the tools we used:

\begin{itemize}
\item Salient object identifier: This is a VLM that takes in the input image and returns the list of important objects in the scene.
\item Segmentation model: This model generates segmentation masks for the list of salient objects identified. 
\item Depth model: We first use a model that given an image estimates depth from camera for each pixel. To get the depth of each salient object, we average the depth values for all pixels belonging to that object. We make use of the segmention masks to identify which pixels belong to an object.
\item Density table: For each image in dataset, we ask a VLM to identify the materials of the salient objects. We arrive at a list of 22 materials and specify their densities in kilograms per $m^3$ to get a table.
\item Labeled angle schemas: We provided a set of images where in each image two lines intersect at a specified angle. The angles we considered for this set are 0, 45, 90, 135, and 180 degrees.
\end{itemize}

Not all of these above tool outputs are relevant to every category of estimation. To avoid overwhelming the context with unnecessary information, we only the following tool outputs for specific categories:

\begin{itemize}
\item Weight: Depth of salient objects plus density table.
\item Angle: Depth of salient objects plus labeled angle schemas.
\item Volume: Only depth of salient objects.
\item Degree of stability: Only depth of salient objects.
\item Geometric fit: Only depth of salient objects.
\end{itemize}

\subsubsection{Intervention 3: Fine-tuning open weight VLM}
We fine-tune small open-weight VLMs on a training dataset based on weight and volume estimation tasks. To derive the dataset, we use automatically selected a subset of images from publicly available image-classification datasets. For inclusion in our training set, an image should be such that it's feasible to pose and answer an estimation question using it. For example, a good candidate image should not have an object in close focus; there should be diverse objects in surroundings which can act as references or anchors during reasoning. To select such images, we use a strong VLM with a prompt that specifies our desidarata. The strong VLM then gives us scores for whether the input image can be a good candidate for one of our five estimation categories. We found that most of the selected images belonged to weight and volume estimation. Since the fine-tuning dataset must at least comprise a few thousand samples, we restricted ourselves to weight and volume estimation fine-tuning to limit compute expense.

We used parameter-efficient fine-tuning using low-rank adaptation (LoRA).
